% Vorlage Schema / Prinzipskizze (TikZ) im Leitfaden-Stil
%
% Uebersetzen mit LuaLaTeX. Im selben Ordner muessen liegen:
%   vfdbstil.sty, vfdbstil-werte.tex
% Ergebnis: PDF in fester Groesse -> in Overleaf nach media/ hochladen und
% ohne Skalierung einbinden: \includegraphics{media/XX-name.pdf}
%
% Koordinaten in mm, Ursprung = Schnittpunkt der Achsen.
% Achsenflaeche: \vfdbB x \vfdbH (Standard 103 x 64 mm).
% SPDX-License-Identifier: MIT
\documentclass[border=0pt]{standalone}
\usepackage[ngerman]{babel}
\usepackage[schriften]{vfdbstil}

\begin{document}
% [Format]{unterer Rand mm}{oberer Rand mm}; mit Phasenbalken unten ca. 14 mm
\begin{vfdbschema}[standard]{14}{3}
  % Flaeche und Kurve
  \fill[vfdb flaeche] (0,0) -- plot[domain=0:40, samples=60] (\x, {40*(\x/40)^2})
                      -- (70, 40) -- (90, 0) -- cycle;
  \draw[vfdb kurve=1] plot[domain=0:40, samples=60] (\x, {40*(\x/40)^2}) -- (70, 40) -- (90, 0);
  \draw[vfdb hilfslinie] (40, 0) -- (40, 40);

  % Achsen als Pfeile
  \vfdbPfeilachsen{Wärmefreisetzungsrate $\dot{Q}$}

  % Hinweis: Pfeil waagerecht oder senkrecht
  \node[anchor=west] (h) at (50, 20) {Beschriftung};
  \draw[vfdb hinweis] (h.west) -- (32, 20);

  % Beschriftungen an den Achsen (8 pt)
  \vfdbXmarke{40}{$t_{1}$}
  \vfdbYmarke{40}{$\dot{Q}_\mathrm{max}$}   % beschreibender Index aufrecht

  % Phasen unter der Zeitachse
  \vfdbPhase{0}{40}{Phase 1}
  \vfdbPhase{40}{90}{Phase 2}
\end{vfdbschema}
\end{document}
